\section{Motivating Example}
\label{sec:motivation}
Consider an \emph{illustrative} skill $S_1$ that wraps database calls in a bounded retry.  It was learned on a repository whose client
has at-least-once semantics; on a target whose client is exactly-once, the retry duplicates writes under partial failure.  The skill
is authorized (it edits application code only), its replay tests pass on the source, and the target's environment fingerprint can
agree on every recorded variable while differing in the one that matters, delivery semantics, which nobody recorded.  Only evidence
gathered on the target can reveal the problem.  Now add a skill $S_2$ that turns the call into a coroutine.  If $S_1$ wraps the call
before it is awaited, its retry ``succeeds'' on the coroutine object and the real failure is never raised: each skill is beneficial
alone, the pair is harmful.  Theorem~\ref{thm:composition} shows that nothing about the two individual risks prevents this, which is
why compositions are cells of their own and why the execution rule refuses to run overlapping skills together until their composite is
admitted.
